\documentclass[11pt]{article}
\usepackage{kotex}
\usepackage{forest}
\usepackage[a4paper,top=2.5cm,bottom=2.5cm,left=2cm,right=2cm]{geometry}
\usepackage{fancyhdr}
\pagestyle{fancy}
\fancyhf{}
\fancyhead[R]{Syntax(001) - quiz(2)}
\fancyfoot[L]{\nouppercase{\leftmark}}
\fancyfoot[R]{\thepage}
\renewcommand{\sectionmark}[1]{\markboth{#1}{}}
\renewcommand{\headrulewidth}{0.4pt}
\renewcommand{\footrulewidth}{0.4pt}
\setlength{\headheight}{14pt}
\fancypagestyle{firstpage}{%
  \fancyhf{}
  \fancyfoot[L]{}
  \fancyfoot[R]{\thepage}
  \renewcommand{\headrulewidth}{0pt}
  \renewcommand{\footrulewidth}{0.4pt}
}
\setcounter{secnumdepth}{0}

\title{2026 Fall Syntax(001): \\ \large Quiz 2}
\author{강초롱 / 14-602}
\date{November 4, 2026}

\begin{document}
\maketitle
\thispagestyle{firstpage}

\section*{CPS1. Hungarian DPs (Chapter 7)}

In the text above, we argued that the structure of genitive constructions in English looks like:

\begin{center}
\begin{forest}
for tree={align=center, s sep=10mm, l sep=8mm}
[DP
  [DP [possessor]]
  [D$'$ 
    [D [\textit{'s}]]
    [NP [N$'$ [N [\textit{possessed}]] [\ldots]]]
  ]
]
\end{forest}
\end{center}

Consider the following data from Hungarian. Does the possessor DP appear in the same place as the English one? Assume the determiners \textit{az} and \textit{a} modify the possessed noun, not the possessor. The ending on the word \textit{kalapom/kalapod/kalapja} varies depending upon the possessor. This does not affect the answer to this question. (Data from Szabolcsi 1994.)

\begin{enumerate}
  \item[a)]
  \begin{tabular}[t]{@{}l@{\quad}l@{\quad}l@{}}
    \textit{az} & \textit{én} & \textit{kalapom}\\
    the & I & hat-1SG
  \end{tabular}\\
  ``my hat''
  \item[b)]
  \begin{tabular}[t]{@{}l@{\quad}l@{\quad}l@{}}
    \textit{a} & \textit{te} & \textit{kalapod}\\
    the & you & hat-2SG
  \end{tabular}\\
  ``your hat''
\end{enumerate}

Hungarian has another possessive construction, seen in (c).

\begin{enumerate}
  \item[c)]
  \begin{tabular}[t]{@{}l@{\quad}l@{\quad}l@{}}
    \textit{Marinak} & \textit{a} & \textit{kalapja}\\
    Mary & the & hat-3SG
  \end{tabular}\\
  ``Mary's hat''
\end{enumerate}

Where is the possessor DP in (c)? Explain your answer.

\newpage
\section*{GPS3 \& CPS1. Irish Finite Complementizers (Chapter 9)}

In English, complementizers are primarily sensitive to whether or not the embedded TP is an infinitive marked with \textit{to}. In Modern Irish, finite complementizers vary along an additional line. Examine the data below. The extra \textit{m} and \textit{h} on either side of \textit{b} in the embedded verb are effects of initial consonant mutation and are irrelevant here. Assume the data are representative of the patterns in the language.

\begin{enumerate}
  \item[(a)]
  \begin{tabular}[t]{@{}l@{\quad}l@{\quad}l@{\quad}l@{\quad}l@{\quad}l@{\quad}l@{}}
    \textit{Ceapaim} & \textit{go} & \textit{mbuaileann} & \textit{sé} & \textit{le} & \textit{Seán} & \textit{inniu.}\\
    think.1SG & that & meets & he & with & John & today
  \end{tabular}\\
  ``I think that he meets with John today.''

  \item[(b)]
  \begin{tabular}[t]{@{}l@{\quad}l@{\quad}l@{\quad}l@{\quad}l@{\quad}l@{\quad}l@{}}
    \textit{Ceapaim} & \textit{gur} & \textit{bhuail} & \textit{sé} & \textit{le} & \textit{Seán} & \textit{inné.}\\
    think.1SG & that & met & he & with & John & yesterday
  \end{tabular}\\
  ``I think that he met with John yesterday.''

  \item[(c)]
  \begin{tabular}[t]{@{}l@{\quad}l@{\quad}l@{\quad}l@{\quad}l@{\quad}l@{\quad}l@{}}
    \textit{Ceapaim} & \textit{go} & \textit{mbuailfidh} & \textit{sé} & \textit{le} & \textit{Seán} & \textit{inniu.}\\
    think.1SG & that & meet.FUT & he & with & John & today
  \end{tabular}\\
  ``I think that he will meet with John today.''

  \item[(d)]
  \begin{tabular}[t]{@{}l@{\quad}l@{\quad}l@{\quad}l@{\quad}l@{\quad}l@{\quad}l@{}}
    \textit{Ceapaim} & \textit{gur} & \textit{bhuailfeadh} & \textit{sé} & \textit{le} & \textit{Seán} & \textit{inné.}\\
    think.1SG & that & meet.COND & he & with & John & yesterday
  \end{tabular}\\
  ``I think that he would meet with John yesterday.''
\end{enumerate}

Taking these data together with the Irish complementizer data in General Problem Set 2, what do the theta grids for \textit{go} and \textit{gur} look like? Propose a theta grid for each complementizer and explain how the data support your analysis.

\newpage
\section*{3. Tree Drawing}

Draw an X-bar theoretic tree for each sentence below. Assume the DP proposal and the projections ProgP, PerfP, and Voice$_{passive}$P have been introduced. Use the appropriate projections in each tree and include all relevant functional structure.

\vspace{2em}
\begin{enumerate}
  \item The tuna will be being eaten.

  \vspace{8cm}

  \item Calvin had been eating.

  \vspace{8cm}
\end{enumerate}
\end{document}








